\documentclass[11pt,a4paper]{article}
\usepackage[margin=1in]{geometry}
\usepackage[T1]{fontenc}
\usepackage[utf8]{inputenc}
\usepackage{lmodern}
\usepackage{hyperref}
\usepackage{enumitem}
\usepackage{xcolor}
\usepackage{titlesec}
\usepackage{parskip}

\hypersetup{
  colorlinks=true,
  linkcolor=black,
  urlcolor=blue
}

\titleformat{\section}{\large\bfseries}{}{0em}{}
\titleformat{\subsection}{\normalsize\bfseries}{}{0em}{}
\setlist[itemize]{leftmargin=1.5em}

\begin{document}

\begin{center}
    {\LARGE \textbf{Kaze Investor Brief}}\\[0.5em]
    {\large Venue Experience, Access, and Commerce Platform}\\[0.75em]
    April 2026\\[0.75em]
    \textbf{Founder:} Oreste Gabo\\
    \href{mailto:orestegabo@icloud.com}{orestegabo@icloud.com} \\
    \href{https://orestegabo.dev}{https://orestegabo.dev}
\end{center}

\section{Executive Summary}

Kaze is a venue experience, access, and commerce platform designed for Rwanda and East Africa. It began as a hotel guest experience product, but its architecture and product direction support a much broader opportunity: reusable venue maps, reservations, access control, local payments, and value-added services for multiple venue categories.

Kaze can serve:
\begin{itemize}
    \item hotels
    \item conference venues
    \item wedding venues
    \item stadiums
    \item government buildings
    \item other structured physical spaces such as airports and transport environments
\end{itemize}

\section{Problem}

Many venue experiences remain fragmented:
\begin{itemize}
    \item navigation is unclear inside large venues
    \item services are requested through phone calls, ad-hoc chats, or staff friction
    \item venue reservations are often poorly digitized
    \item local payment options are not always integrated cleanly
    \item venues lose revenue when customers bypass the platform and negotiate directly
    \item reusable spatial data is rarely treated as a product asset across multiple apps
\end{itemize}

\section{Solution}

Kaze combines:
\begin{itemize}
    \item venue maps and indoor wayfinding
    \item digital access and Kaze Pass entry control
    \item guest and attendee experience flows
    \item venue reservations
    \item local payment support for Rwanda
    \item add-on services that increase booking value
\end{itemize}

This creates a platform rather than a single-purpose app.

\section{Market Direction}

Kaze does not need to begin with hotel room distribution, where large international players already dominate. Instead, it can begin where local advantage is stronger:
\begin{itemize}
    \item conference room reservations
    \item wedding venue reservations
    \item event access and pass-controlled entry
    \item venue-specific add-ons
    \item locally relevant payments
\end{itemize}

\section{Revenue Opportunities}

Potential revenue sources:
\begin{itemize}
    \item SaaS subscriptions for venue operators
    \item booking commissions
    \item add-on service commissions
    \item setup and import fees for venue map digitization
    \item white-label branded venue apps
    \item enterprise or government deployments
    \item premium analytics and operational dashboards
\end{itemize}

\section{Local Payment Advantage}

For Rwanda, Kaze can support locally understood payment methods such as:
\begin{itemize}
    \item MTN MoMo
    \item Airtel Money
    \item BK / RSwitch compatible payment flows
    \item other locally relevant mobile and bank-backed rails
\end{itemize}

This is important because local payment fit can improve conversion and operator trust.

\section{Platform Defensibility}

Customers may try to bypass a booking platform and negotiate directly with a hotel or venue. Kaze can reduce this risk by making the platform more valuable than the bypass.

Examples of in-platform value:
\begin{itemize}
    \item layout planning for conference and wedding venues
    \item Kaze Pass gated entry after payment or reservation approval
    \item event styling and decor add-ons
    \item cleaning services
    \item insurance add-ons
    \item photography, videography, and live streaming services
\end{itemize}

\section{Why Kaze Can Expand Beyond Hotels}

The same underlying platform model can support multiple physical-space businesses:
\begin{itemize}
    \item maps
    \item access rules
    \item space numbering
    \item reservations
    \item payments
    \item pass-controlled entry
\end{itemize}

This gives Kaze a chance to become infrastructure for venue software, not just one app.

\section{Current State}

The product already has:
\begin{itemize}
    \item Kotlin Multiplatform app structure
    \item Android and iPhone app targets
    \item indoor map UI direction
    \item service-request flows
    \item event and explore flows
    \item digital access pass UI
    \item initial backend and shared architecture foundation
\end{itemize}

\section{Near-Term Focus}

\begin{enumerate}
    \item Finalize the reusable venue-map data model.
    \item Build the first conference and wedding reservation flow.
    \item Add Rwanda payment integrations.
    \item Introduce layout planning for chairs, tables, and attendee-based configurations.
    \item Add Kaze Pass-gated event entry after reservation or payment.
    \item Launch pilot partnerships with local venues.
\end{enumerate}

\section{Contact}

For investor, pilot, or partnership discussions:

\begin{itemize}
    \item \textbf{Email:} \href{mailto:orestegabo@icloud.com}{orestegabo@icloud.com}
    \item \textbf{Website:} \href{https://orestegabo.dev}{https://orestegabo.dev}
\end{itemize}

\end{document}
